\documentclass[onecolumn,10pt,titlepage]{article}
\usepackage[a4paper,top=2.5cm,bottom=2cm,left=2cm,right=2cm,marginparwidth=1.75cm,headheight=28pt]{geometry}
\usepackage[utf8]{inputenc}
\usepackage[spanish]{babel} 
\usepackage{natbib}
\usepackage{multicol}
\usepackage{mathastext} %BORRAR O DEJAR: Cambia estilo de la fuente matemática y la deja con aspecto mas "técnico".
\usepackage{graphicx}
\usepackage{hyperref}
\hypersetup{
    colorlinks,
    citecolor=black,
    filecolor=black,
    linkcolor=black,
    urlcolor=black
}

%Helvetica
\renewcommand{\familydefault}{\sfdefault}
\usepackage[scaled=1]{helvet}
\usepackage[format=plain,
            labelfont={bf,it},
            textfont=it]{caption}
%--------------------------------------




\newcommand{\glossentry}[2]{$#1$ \indent #2 \par \vspace{.4cm} }

\title{Informe Técnico - ITBA}

\author{Patricio Whittingslow}
%========================> Comienza Documento
\begin{document}

\begin{titlepage}
\centering

{ \large Instituto Tecnológico de Buenos Aires  \par }
\vspace{2cm}
{\Large \scshape Física V - 93.45 \par}
\vspace{2cm}
{\Huge \scshape Explicación del efecto Hall Cuántico \par }
\vspace{1cm}
{\large \bf Grupo 10 \par}
\vspace{0cm}
\textsc{\large Autor -- Legajo \\ Coautor -- Legajo}
\vspace{2cm}
{\par \large Fecha de realización: \today \par}
\vspace{1cm}
{\large Fecha de entrega: .......................................\par}
\vspace{2.5cm}
{\large Firma del docente: .......................................}
\vspace{3cm}
\begin{figure}[htb!]
\centering
\includegraphics[width=6cm]{logoitba.png}
\end{figure}
\end{titlepage}


{\textbf{Problema o Resumen}}\par
Un resumen del problema a tratar debería tener menos de 200 palabras y no incluir referencias. Un resumen del problema a tratar debería tener menos de 200 palabras y no incluir referencias. Un resumen del problema a tratar debería tener menos de 200 palabras y no incluir referencias. Un resumen del problema a tratar debería tener menos de 200 palabras y no incluir referencias. Un resumen del problema a tratar debería tener menos de 200 palabras y no incluir referencias.

\begin{multicols}{2}
\section*{Glosario}
\glossentry{\pi}{Valor de pi.}
\glossentry{g}{Aceleración de la gravedad en la tierra.}
%\glossentry{x}{}
\end{multicols}


\tableofcontents


\section{Introducción Teórica}
Para que vean como se ve la matemática con el paquete \verb|mathastext|:
$e^{i\pi}+1=0$ 

$$\int^\infty_0 \frac{1}{x^2}.dx $$

Esto se puede cambiar simplemente borrando la linea de código \verb|7|, por defecto.

\subsection*{Numeración}
En general no se numera las subsecciones. Esto se logra agregando un asterisco: \verb|\subsection*{titulo}|
\subsection*{Referencias}
Citar es facil:
\citet*[pag. 4]{silverstein1964giving}
y \cite{serway2010fisica}, \citep{Erdos65}. Agreguen sus referencias a \verb|labibliografia.bib|.

\subsubsection*{Minisección}
Se puede ir hasta un nivel bajo.
\section{Método}
Otros deberían poder repetir el trabajo con la información disponible aquí.
\section{Resultados}
Hagan click sobre el indice o la referencias para navegar el documento.
\section{Discusión o conclusión}
La discusión debería ser corta y no tener subsecciones.

\clearpage 


\bibliography{labibliografia} % Indica archivo
\bibliographystyle{plainnat} 

\end{document}